\chapter{Risultati e analisi}
\label{ch:risultati}

Questo capitolo raccoglie le misure dei test e le analizza. Salvo diverso avviso i numeri sono medie su tre esecuzioni
a freddo, ovvero la cache viene svuotata prima di ogni esecuzione, per lo scaling per nodi, e misure a caldo per lo scaling sui dati. 
Il dataset di partenza è delle stesse dimensioni su entrambi i sistemi (30 ditte, 554 dipendenti), generato con lo script dedicato.

\section{Baseline in lettura e i due regimi}
\label{sec:baseline}

Su Citus a un nodo, sotto carico sostenuto, emergono due regimi netti. Le query \textbf{single-shard} (cedolino singolo, cedolini di una ditta, riepiloghi per un tenant) stanno a \textbf{1--3\,ms} con throughput
fino a ${\sim}680$ op/s: colpiscono un solo shard e il coordinator ha solo un overhead fisso di instradamento. Le analitiche \textbf{cross-shard} (percentili, aggregazioni cross-tenant, spesa per categoria)
stanno a \textbf{55--85\,ms} con 13--18 op/s: ogni worker satura la CPU per lo scan (100--190\,\%, multi-core)
e il coordinator (${\sim}40\,\%$) fa il merge. Questa distinzione tra i due regimi si ritrova in tutti gli esperimenti seguenti.

\section{Scalabilità per numero di nodi (Citus)}
\label{sec:scaling-nodi}

A parità di dati si passa da 1 a 4 nodi, ridistribuendo gli shard a ogni passo. Le query cross-shard scalano
in modo \textbf{quasi lineare} (\Cref{tab:scaling-nodi,fig:scaling-nodi}): con 4 nodi lo speedup è
\textbf{3.8--4.6$\times$}. Le single-shard restano invece piatte: colpiscono un solo shard su un solo nodo,
quindi aggiungere worker non le tocca (ma sono già velocissime). Il lavoro si distribuisce: la CPU e il tempo
di esecuzione \emph{per worker} calano proporzionalmente (per L04, la CPU del worker passa da 142\,\% a 24\,\%
e le righe per worker si dimezzano a ogni raddoppio).

\begin{table}[htbp]
  \centering
  \caption{Latenza (ms) al crescere dei nodi e speedup 1$\to$4. Cross-shard vs single-shard.}
  \label{tab:scaling-nodi}
  \begin{tabular}{lrrrrr}
    \toprule
    \textbf{Query} & \textbf{1 nodo} & \textbf{2 nodi} & \textbf{3 nodi} & \textbf{4 nodi} & \textbf{speedup} \\
    \midrule
    L04 gender gap mediano & 71.9 & 45.0 & 23.2 & 18.3 & \textbf{3.9$\times$} \\
    L05 top-N retribuzioni & 65.8 & 33.7 & 19.4 & 14.9 & \textbf{4.4$\times$} \\
    L06 assenteismo/mese   & 61.4 & 31.2 & 19.9 & 15.4 & \textbf{4.0$\times$} \\
    L08 trend costo lavoro & 56.1 & 24.0 & 15.9 & 12.1 & \textbf{4.6$\times$} \\
    L11 spesa/categoria    & 75.6 & 52.7 & 25.3 & 20.1 & \textbf{3.8$\times$} \\
    L12 ferie residue      & 75.4 & 59.7 & 25.2 & 19.5 & \textbf{3.9$\times$} \\
    \midrule
    L01 cedolino singolo \emph{(single-shard)} & 1.5 & 1.5 & 1.5 & 1.5 & 1.0$\times$ \\
    \bottomrule
  \end{tabular}
\end{table}

Il dettaglio completo di \emph{tutte} le metriche raccolte (latenza, throughput, CPU per worker, tempo di
esecuzione e righe elaborate per worker) per ognuna delle 13 query è in \Cref{tab:scaling-nodi-full}: mostra
esplicitamente come, all'aumentare dei nodi, il lavoro \emph{per worker} (CPU, \texttt{exec}, righe) si
riduca, mentre l'overhead sul coordinator resta pressoché costante.

\begingroup
\footnotesize
\begin{longtable}{@{}llrrrr@{}}
\caption{Scalabilità per numero di nodi (Citus): tutte le metriche raccolte, per le 13 query, a 1--4 nodi. Le metriche per-worker sono la media sui worker attivi.}\label{tab:scaling-nodi-full}\\
\toprule \textbf{Query} & \textbf{Parametro} & \textbf{1 nodo} & \textbf{2 nodi} & \textbf{3 nodi} & \textbf{4 nodi} \\ \midrule \endfirsthead
\toprule \textbf{Query} & \textbf{Parametro} & \textbf{1 nodo} & \textbf{2 nodi} & \textbf{3 nodi} & \textbf{4 nodi} \\ \midrule \endhead
\multirow{5}{*}{L01 cedolino singolo} & latenza (ms)     & 1.5 & 1.5 & 1.5 & 1.5 \\
 & throughput (op/s) & 680 & 666 & 669 & 688 \\
 & CPU worker (\%)   & 11 & 6 & 4 & 3 \\
 & exec worker (ms)  & 138 & 63 & 47 & 34 \\
 & righe/worker      & 6797 & 3331 & 2230 & 1718 \\ \midrule
\multirow{5}{*}{L02 cedolini ditta/periodo} & latenza (ms)     & 1.5 & 1.6 & 1.6 & 1.5 \\
 & throughput (op/s) & 669 & 636 & 638 & 654 \\
 & CPU worker (\%)   & 15 & 7 & 5 & 4 \\
 & exec worker (ms)  & 171 & 82 & 56 & 43 \\
 & righe/worker      & 6699 & 3182 & 2127 & 1636 \\ \midrule
\multirow{5}{*}{L03 costo per centro} & latenza (ms)     & 2.3 & 2.3 & 2.3 & 2.3 \\
 & throughput (op/s) & 438 & 431 & 437 & 433 \\
 & CPU worker (\%)   & 27 & 14 & 9 & 7 \\
 & exec worker (ms)  & 646 & 315 & 213 & 161 \\
 & righe/worker      & 210192 & 103456 & 69909 & 51988 \\ \midrule
\multirow{5}{*}{L04 gender gap mediano} & latenza (ms)     & 71.9 & 45.0 & 23.2 & 18.3 \\
 & throughput (op/s) & 14 & 23 & 43 & 55 \\
 & CPU worker (\%)   & 142 & 57 & 26 & 24 \\
 & exec worker (ms)  & 531 & 335 & 367 & 347 \\
 & righe/worker      & 77375 & 63064 & 79899 & 75667 \\ \midrule
\multirow{5}{*}{L05 top-N retribuzioni} & latenza (ms)     & 65.8 & 33.7 & 19.4 & 14.9 \\
 & throughput (op/s) & 15 & 31 & 52 & 67 \\
 & CPU worker (\%)   & 129 & 33 & 20 & 20 \\
 & exec worker (ms)  & 327 & 237 & 244 & 266 \\
 & righe/worker      & 54840 & 55080 & 61960 & 60240 \\ \midrule
\multirow{5}{*}{L06 assenteismo per mese} & latenza (ms)     & 61.4 & 31.2 & 19.9 & 15.4 \\
 & throughput (op/s) & 16 & 32 & 50 & 65 \\
 & CPU worker (\%)   & 122 & 30 & 21 & 17 \\
 & exec worker (ms)  & 1700 & 1225 & 1224 & 923 \\
 & righe/worker      & 245000 & 242750 & 251667 & 187125 \\ \midrule
\multirow{5}{*}{L07 straordinari per centro} & latenza (ms)     & 3.1 & 3.1 & 3.1 & 3.2 \\
 & throughput (op/s) & 321 & 322 & 326 & 317 \\
 & CPU worker (\%)   & 41 & 21 & 14 & 11 \\
 & exec worker (ms)  & 1631 & 889 & 603 & 441 \\
 & righe/worker      & 12843 & 6435 & 4343 & 3165 \\ \midrule
\multirow{5}{*}{L08 trend costo lavoro} & latenza (ms)     & 56.1 & 24.0 & 15.9 & 12.1 \\
 & throughput (op/s) & 18 & 42 & 63 & 82 \\
 & CPU worker (\%)   & 106 & 14 & 12 & 12 \\
 & exec worker (ms)  & 466 & 368 & 360 & 351 \\
 & righe/worker      & 42880 & 50080 & 50267 & 49500 \\ \midrule
\multirow{5}{*}{L09 organico per livello} & latenza (ms)     & 57.7 & 24.5 & 17.0 & 12.8 \\
 & throughput (op/s) & 17 & 41 & 59 & 78 \\
 & CPU worker (\%)   & 106 & 15 & 14 & 13 \\
 & exec worker (ms)  & 186 & 134 & 127 & 125 \\
 & righe/worker      & 51852 & 60991 & 58640 & 58234 \\ \midrule
\multirow{5}{*}{L10 cedolino completo} & latenza (ms)     & 82.7 & 63.3 & 27.9 & 21.7 \\
 & throughput (op/s) & 12 & 16 & 36 & 46 \\
 & CPU worker (\%)   & 189 & 91 & 33 & 31 \\
 & exec worker (ms)  & 169 & 97 & 116 & 112 \\
 & righe/worker      & 121 & 79 & 120 & 115 \\ \midrule
\multirow{5}{*}{L11 spesa per categoria} & latenza (ms)     & 75.6 & 52.7 & 25.3 & 20.1 \\
 & throughput (op/s) & 13 & 19 & 40 & 50 \\
 & CPU worker (\%)   & 154 & 68 & 27 & 25 \\
 & exec worker (ms)  & 1464 & 535 & 524 & 543 \\
 & righe/worker      & 73497 & 52907 & 73066 & 69158 \\ \midrule
\multirow{5}{*}{L12 ferie residue} & latenza (ms)     & 75.4 & 59.7 & 25.2 & 19.5 \\
 & throughput (op/s) & 13 & 17 & 40 & 51 \\
 & CPU worker (\%)   & 165 & 82 & 32 & 30 \\
 & exec worker (ms)  & 909 & 439 & 545 & 528 \\
 & righe/worker      & 3990 & 2520 & 3963 & 3842 \\ \midrule
\multirow{5}{*}{L13 riepilogo assenze} & latenza (ms)     & 3.4 & 3.5 & 3.5 & 3.4 \\
 & throughput (op/s) & 291 & 290 & 289 & 291 \\
 & CPU worker (\%)   & 38 & 19 & 12 & 10 \\
 & exec worker (ms)  & 923 & 453 & 305 & 230 \\
 & righe/worker      & 29320 & 14487 & 9628 & 7281 \\ \midrule
\bottomrule
\end{longtable}
\endgroup


\begin{figure}[htbp]
  \centering
  \includegraphics[width=0.82\textwidth]{fig-scaling-nodi-citus}
  \caption{Citus, scaling per nodi (tutte le 13 letture, scala log; codice query in coda a ogni linea): le
    cross-shard calano nettamente al crescere dei nodi, le single-shard restano piatte.}
  \label{fig:scaling-nodi}
\end{figure}

\begin{lstlisting}[style=sobrio, language=bash, caption={Esecuzione: aggiunta di un worker e rebalance, poi 3 run a freddo delle 13 letture e media (ripetuto a 1, 2, 3, 4 nodi).}]
bash citus-add-node.sh 10.0.1.5 && bash citus-rebalance.sh       # ad ogni passo di scaling
for i in 1 2 3; do bash azzera-cache-citus.sh; \
  bash esegui-suite-citus.sh 10 queries/relational/letture/*.sql; done
bash media-citus.sh
\end{lstlisting}

\section{Scalabilità per volume di dati (Citus)}
\label{sec:scaling-dati}

A numero di nodi fisso (4) si aumenta il volume dei dati. Le cross-shard \textbf{degradano in modo
super-lineare} (\Cref{tab:scaling-dati,fig:scaling-dati}): da 30 a 480 ditte L11 passa da 20 a 114\,ms
(5.7$\times$), con un salto netto tra 240 e 480 ditte. Le single-shard restano \textbf{piatte} (${\sim}1.5$\,ms
a ogni scala): indicizzate e mirate a un tenant, sono insensibili al volume globale, che è il vantaggio del
multitenant co-locato. Le letture da disco (\texttt{blk\_read}) restano nulle: a questa scala i dati stanno
ancora in RAM, quindi il degrado è \emph{CPU-bound} (più righe da scandire in cache). L'esperimento si è
fermato a 480 ditte perché a quel punto la tendenza era già chiara e il carico risultava distribuito su tutti
i nodi: proseguire non avrebbe aggiunto informazione.

\begin{table}[htbp]
  \centering
  \caption{Latenza (ms) al crescere dei dati, a 4 nodi fissi.}
  \label{tab:scaling-dati}
  \begin{tabular}{lrrrr}
    \toprule
    \textbf{Ditte} & \textbf{L01 single-shard} & \textbf{L04 gender gap} & \textbf{L06 assenteismo} & \textbf{L11 spesa} \\
    \midrule
    30 (base) & 1.5 & 18.3 & 15.4 & 20.1 \\
    60        & 1.5 & 20.4 & 17.5 & 24.6 \\
    120       & 1.5 & 24.0 & 25.1 & 34.9 \\
    240       & 1.4 & 30.9 & 51.2 & 66.0 \\
    480       & 1.6 & \textbf{68.3} & \textbf{73.1} & \textbf{114.1} \\
    \bottomrule
  \end{tabular}
\end{table}

Il dettaglio completo di tutte le metriche per le 13 query, a ogni scala, è in \Cref{tab:scaling-dati-full}:
si vede come il tempo di esecuzione e le righe elaborate per worker crescano col volume, mentre le query
single-shard restano invariate.

\begingroup
\footnotesize
\begin{longtable}{@{}llrrrrr@{}}
\caption{Scalabilità per volume di dati (Citus, 4 nodi fissi): tutte le metriche per le 13 query, da 30 a 480 ditte.}\label{tab:scaling-dati-full}\\
\toprule \textbf{Query} & \textbf{Parametro} & \textbf{30} & \textbf{60} & \textbf{120} & \textbf{240} & \textbf{480} \\* \midrule \endfirsthead
\toprule \textbf{Query} & \textbf{Parametro} & \textbf{30} & \textbf{60} & \textbf{120} & \textbf{240} & \textbf{480} \\* \midrule \endhead
\multirow{5}{*}{L01 cedolino singolo} & latenza (ms)     & 1.5 & 1.5 & 1.5 & 1.4 & 1.6 \\
 & throughput (op/s) & 688 & 659 & 655 & 705 & 643 \\
 & CPU worker (\%)   & 3 & 6 & 6 & 5 & 4 \\
 & exec worker (ms)  & 34 & 36 & 36 & 38 & 34 \\
 & righe/worker      & 1718 & 1647 & 1636 & 1762 & 1607 \\ \midrule
\multirow{5}{*}{L02 cedolini ditta/periodo} & latenza (ms)     & 1.5 & 1.5 & 1.5 & 1.6 & 1.5 \\
 & throughput (op/s) & 654 & 659 & 674 & 635 & 682 \\
 & CPU worker (\%)   & 4 & 5 & 4 & 4 & 4 \\
 & exec worker (ms)  & 43 & 44 & 45 & 33 & 36 \\
 & righe/worker      & 1636 & 1648 & 1684 & 1588 & 1706 \\ \midrule
\multirow{5}{*}{L03 costo per centro} & latenza (ms)     & 2.3 & 2.3 & 2.4 & 2.4 & 2.4 \\
 & throughput (op/s) & 433 & 438 & 424 & 414 & 423 \\
 & CPU worker (\%)   & 7 & 7 & 7 & 8 & 8 \\
 & exec worker (ms)  & 161 & 164 & 171 & 175 & 185 \\
 & righe/worker      & 51988 & 52548 & 50820 & 49680 & 50748 \\ \midrule
\multirow{5}{*}{L04 gender gap mediano} & latenza (ms)     & 18.3 & 20.4 & 24.0 & 30.9 & 68.3 \\
 & throughput (op/s) & 55 & 49 & 42 & 32 & 15 \\
 & CPU worker (\%)   & 24 & 26 & 26 & 28 & 38 \\
 & exec worker (ms)  & 347 & 623 & 947 & 1430 & 1325 \\
 & righe/worker      & 75667 & 164762 & 283664 & 454734 & 411049 \\ \midrule
\multirow{5}{*}{L05 top-N retribuzioni} & latenza (ms)     & 14.9 & 14.8 & 15.8 & 15.7 & 17.4 \\
 & throughput (op/s) & 67 & 68 & 63 & 64 & 58 \\
 & CPU worker (\%)   & 20 & 22 & 23 & 27 & 30 \\
 & exec worker (ms)  & 266 & 423 & 587 & 953 & 1475 \\
 & righe/worker      & 60240 & 89401 & 98425 & 102080 & 92320 \\ \midrule
\multirow{5}{*}{L06 assenteismo per mese} & latenza (ms)     & 15.4 & 17.5 & 25.1 & 51.2 & 73.1 \\
 & throughput (op/s) & 65 & 57 & 40 & 20 & 14 \\
 & CPU worker (\%)   & 17 & 32 & 39 & 68 & 91 \\
 & exec worker (ms)  & 923 & 2445 & 3337 & 3655 & 5274 \\
 & righe/worker      & 187125 & 326326 & 262342 & 134554 & 94256 \\ \midrule
\multirow{5}{*}{L07 straordinari per centro} & latenza (ms)     & 3.2 & 3.2 & 3.2 & 3.2 & 3.3 \\
 & throughput (op/s) & 317 & 317 & 309 & 312 & 307 \\
 & CPU worker (\%)   & 11 & 11 & 11 & 11 & 12 \\
 & exec worker (ms)  & 441 & 478 & 474 & 483 & 480 \\
 & righe/worker      & 3165 & 3173 & 3093 & 3122 & 3068 \\ \midrule
\multirow{5}{*}{L08 trend costo lavoro} & latenza (ms)     & 12.1 & 12.6 & 13.9 & 14.9 & 17.8 \\
 & throughput (op/s) & 82 & 79 & 72 & 67 & 56 \\
 & CPU worker (\%)   & 12 & 14 & 18 & 27 & 39 \\
 & exec worker (ms)  & 351 & 673 & 1148 & 2091 & 3460 \\
 & righe/worker      & 49500 & 68904 & 66960 & 64608 & 54144 \\ \midrule
\multirow{5}{*}{L09 organico per livello} & latenza (ms)     & 12.8 & 13.5 & 14.6 & 16.8 & 20.2 \\
 & throughput (op/s) & 78 & 74 & 69 & 60 & 49 \\
 & CPU worker (\%)   & 13 & 13 & 14 & 15 & 17 \\
 & exec worker (ms)  & 125 & 210 & 336 & 546 & 881 \\
 & righe/worker      & 58234 & 104710 & 185220 & 321186 & 535095 \\ \midrule
\multirow{5}{*}{L10 cedolino completo} & latenza (ms)     & 21.7 & 21.7 & 21.8 & 21.3 & 21.7 \\
 & throughput (op/s) & 46 & 46 & 46 & 47 & 46 \\
 & CPU worker (\%)   & 31 & 32 & 31 & 31 & 30 \\
 & exec worker (ms)  & 112 & 115 & 115 & 119 & 116 \\
 & righe/worker      & 115 & 115 & 115 & 118 & 115 \\ \midrule
\multirow{5}{*}{L11 spesa per categoria} & latenza (ms)     & 20.1 & 24.6 & 34.9 & 66.0 & 114.1 \\
 & throughput (op/s) & 50 & 41 & 29 & 15 & 9 \\
 & CPU worker (\%)   & 25 & 27 & 29 & 42 & 52 \\
 & exec worker (ms)  & 543 & 1102 & 1292 & 1404 & 1718 \\
 & righe/worker      & 69158 & 136854 & 195912 & 213332 & 246070 \\ \midrule
\multirow{5}{*}{L12 ferie residue} & latenza (ms)     & 19.5 & 20.4 & 23.5 & 28.3 & 49.6 \\
 & throughput (op/s) & 51 & 49 & 43 & 35 & 20 \\
 & CPU worker (\%)   & 30 & 35 & 38 & 44 & 65 \\
 & exec worker (ms)  & 528 & 1007 & 1576 & 2569 & 3100 \\
 & righe/worker      & 3842 & 7365 & 12780 & 21240 & 24240 \\ \midrule
\multirow{5}{*}{L13 riepilogo assenze} & latenza (ms)     & 3.4 & 3.6 & 3.6 & 3.6 & 3.5 \\
 & throughput (op/s) & 291 & 279 & 276 & 280 & 289 \\
 & CPU worker (\%)   & 10 & 10 & 10 & 10 & 10 \\
 & exec worker (ms)  & 230 & 233 & 224 & 228 & 193 \\
 & righe/worker      & 7281 & 6978 & 6895 & 7010 & 7220 \\ \midrule
\bottomrule
\end{longtable}
\endgroup


\begin{figure}[htbp]
  \centering
  \includegraphics[width=0.82\textwidth]{fig-scaling-dati-citus}
  \caption{Citus, scaling per dati (tutte le 13 letture, 4 nodi fissi): degrado super-lineare delle
    cross-shard; single-shard insensibili al volume.}
  \label{fig:scaling-dati}
\end{figure}

\begin{lstlisting}[style=sobrio, language=bash, caption={Esecuzione: a 4 nodi fissi, un lotto per volta genera nuovi tenant, carica e misura a caldo (ripetuto fino a 480 ditte).}]
bash scala-dati-citus.sh 31 30 60ditte      # +30 tenant -> 60 ditte; poi 61 60 120ditte; ...
\end{lstlisting}

\section{Scritture distribuite (Citus)}
\label{sec:scritture-citus}

Le scritture, misurate a 4 nodi (\Cref{tab:scritture-citus}), mostrano l'effetto della \textbf{co-locazione}:
\textbf{cinque scritture su sei colpiscono un solo nodo} (quello dello shard del tenant), mentre gli altri tre
worker restano inattivi. La scrittura di un singolo tenant \emph{non si distribuisce} aggiungendo nodi; il
throughput aggregato scala solo distribuendo scritture su tenant diversi. Il costo cresce con le righe/tabelle
toccate e il WAL prodotto: dalla semplice \texttt{UPDATE} (U02, 605 op/s) alla \textbf{transazione}
multi-tabella I03 (17 op/s, ${\sim}1$\,MB di WAL) c'è un fattore ${\sim}35\times$. L'unica scrittura che tocca
\textbf{tutti i nodi} è U03, l'aggiornamento di un dato di riferimento su una reference table (replicata
ovunque): una riga logica, ma applicata su ogni nodo.

\begin{table}[htbp]
  \centering
  \caption{Scritture a 4 nodi (N${=}300$ operazioni per query).}
  \label{tab:scritture-citus}
  \begin{tabular}{llrrl}
    \toprule
    \textbf{Op.} & \textbf{Descrizione} & \textbf{lat (ms/op)} & \textbf{tps} & \textbf{dove scrive} \\
    \midrule
    U02 & licenziamento (update 1 riga)     & 1.7  & 605 & solo 1 nodo \\
    U01 & coordinate bancarie               & 3.7  & 273 & solo 1 nodo \\
    U04 & update concorrente (riga cond.)   & 3.8  & 266 & solo 1 nodo \\
    I01 & inserimento timbratura            & 10.6 & 94  & solo 1 nodo \\
    U03 & modifica dato di riferimento      & 21.1 & 47  & \textbf{tutti i nodi} \\
    I03 & elaborazione cedolino (transazione)& 57.4 & 17  & solo 1 nodo \\
    \bottomrule
  \end{tabular}
\end{table}

\begin{lstlisting}[style=sobrio, language=bash, caption={Esecuzione delle scritture (Citus): N${=}300$ ripetizioni per query, tenant fisso.}]
bash esegui-scritture-citus.sh 300 queries/relational/scritture/*.sql
\end{lstlisting}

\section{Comportamento in presenza di guasti (Citus)}
\label{sec:guasti-citus}

Il comportamento \textbf{CP} è stato messo alla prova eseguendo una query cross-shard (L04) in un ciclo mentre
un worker scelto a caso viene spento e poi riacceso (11 prove). Il risultato è netto e coerente: spegnendo
\emph{qualsiasi} worker, la query fallisce per \textbf{tutta} la finestra di down (15\,s), perché ha bisogno
degli shard di ogni nodo, quindi manca sempre qualcosa e \textbf{non restituisce risultati parziali} (le righe
crollano da $N$ a $0$). L'error-rate medio è ${\sim}51\,\%$, il recupero è \textbf{automatico e rapido}
(prima risposta ${\sim}0.5$\,s dopo il riavvio del container). I quattro worker danno lo stesso error-rate:
ognuno è essenziale, nessuno ridondante (copia singola, \texttt{shard\_replication\_factor}${=}1$). È la
consistenza scelta al posto della disponibilità.

\begin{lstlisting}[style=sobrio, language=bash, caption={Esecuzione del test di guasto (Citus): L04 in loop mentre un worker a caso viene spento (10--25\,s) e riacceso, su 45\,s totali.}]
bash test-guasto-citus.sh queries/relational/letture/L04_gender_gap_mediano.sql random 10 15 45
\end{lstlisting}

\section{MongoDB}
\label{sec:mongo}

\subsection{Scalabilità per shard: il balancer non distribuisce a scala piccola}
Sul cluster documentale il regime ``più shard, stessi dati'' ha rivelato un comportamento che è esso stesso il
risultato più interessante. Con il dataset base (\texttt{cedolini} ${\approx}42$\,MB), il balancer di MongoDB
\textbf{non distribuisce}: con la dimensione di range di default (128\,MB) la soglia di migrazione è 384\,MB,
mai raggiunta, quindi resta un solo chunk su un solo shard. Anche forzando le condizioni, con il range
ridotto a 4\,MB (soglia 12\,MB) e uno \emph{split manuale} in 10 chunk, il balancer, pur attivo (\texttt{mode: full},
888 round eseguiti), \textbf{non ha migrato nulla}. Di conseguenza il secondo shard resta inattivo (0.2\,\%
CPU) e le prestazioni a 1 e 2 shard sono \textbf{identiche} (\Cref{tab:mongo-shard}): aggiungere shard non
porta alcun beneficio. È l'opposto di Citus, dove il rebalance esplicito distribuisce sempre. Il balancer
distribuisce solo a \emph{volume grande}: già a ${\approx}230$ ditte (cedolini ${>}300$\,MB) inizia a spargere
i chunk sugli shard, e il comportamento dipende dunque dalla scala.

\begin{table}[htbp]
  \centering
  \caption{MongoDB, latenza (ms) a 1 e 2 shard: identica, perché il 2\textsuperscript{o} shard resta inattivo.}
  \label{tab:mongo-shard}
  \begin{tabular}{lrr}
    \toprule
    \textbf{Query} & \textbf{1 shard} & \textbf{2 shard} \\
    \midrule
    L04 gender gap (aggregazione) & 7.2  & 7.2  \\
    L11 spesa per categoria        & 12.4 & 12.5 \\
    L01 cedolino singolo (mirata)  & 2.3  & 2.3  \\
    \bottomrule
  \end{tabular}
\end{table}

\begin{lstlisting}[style=sobrio, language=bash, caption={Esecuzione dello scaling per shard (MongoDB): aggiunta automatica degli shard e 3 run a freddo delle letture per passo.}]
export MONGO_PASSWORD=...
bash scala-nodi-mongo.sh 4        # dal numero di shard attuale fino a 4
\end{lstlisting}

\subsection{Scalabilità per volume di dati (MongoDB)}
\label{sec:mongo-dati}
Applicando lo \emph{stesso} protocollo delle controparti Citus (identica suite, misura a caldo, carico
randomizzato, contatori per shard), a numero di shard fisso (4) si aumenta il volume dei dati. Qui si completa
il quadro emerso con lo scaling per shard: mentre a scala piccola il balancer non distribuiva, \textbf{a volume
grande distribuisce davvero}. A ogni scala i
\texttt{cedolini} risultano spartiti su \textbf{tutti e 4 gli shard}, attivi al 48--65\,\% di CPU. Le query
cross-shard degradano super-linearmente (\Cref{tab:mongo-dati}): L04 26$\to$185\,ms, L11 34$\to$308\,ms; le
più pesanti, L06 e L12, crescono a dismisura (302$\to$2834\,ms e 476$\to$4786\,ms) perché aggregano tutti i cedolini
scorrendone gli array annidati: a scala grande il costo dell'embedding sull'intera collezione diventa
dominante. Le query mirate restano piatte (L01 ${\sim}2.6$\,ms).

\begin{table}[htbp]
  \centering
  \caption{MongoDB, scalabilità per dati (4 shard fissi): latenza (ms) al crescere del volume.}
  \label{tab:mongo-dati}
  \begin{tabular}{rrrrrr}
    \toprule
    \textbf{Ditte} & \textbf{L01} & \textbf{L04} & \textbf{L11} & \textbf{L06} & \textbf{L12} \\
    \midrule
     230 & 2.6 & 26.0  & 34.2  & 301.6  & 476.2  \\
     630 & 2.6 & 61.6  & 80.0  & 729.1  & 1174.1 \\
    1430 & 2.7 & 122.2 & 178.0 & 1684.9 & 2759.5 \\
    2630 & 2.7 & 185.2 & 308.0 & 2833.8 & 4786.3 \\
    \bottomrule
  \end{tabular}
\end{table}

Il dettaglio completo per tutte le 13 query e tutte le metriche (latenza, throughput, CPU media per shard,
I/O di disco letto e scritto) è in \Cref{tab:scaling-dati-mongo-full}. Due osservazioni: la CPU cresce su
\emph{tutti} e quattro gli shard (a conferma che il balancer, a questo volume, ha distribuito), e le letture
da disco restano nulle anche a 2630 ditte: il working set è ancora servito dalla cache WiredTiger, quindi
il degrado è dominato dal numero di documenti esaminati, non dall'I/O.

\begingroup
\footnotesize
\begin{longtable}{@{}llrrrrr@{}}
\caption{Scalabilità per volume di dati (MongoDB, 4 shard fissi): tutte le metriche raccolte per le 13 query, da 230 a 2630 ditte. Le metriche per-shard sono la media/somma sui 4 shard.}\label{tab:scaling-dati-mongo-full}\\
\toprule \textbf{Query} & \textbf{Parametro} & \textbf{230} & \textbf{630} & \textbf{1430} & \textbf{2630} \\* \midrule \endfirsthead
\toprule \textbf{Query} & \textbf{Parametro} & \textbf{230} & \textbf{630} & \textbf{1430} & \textbf{2630} \\* \midrule \endhead
\multirow{5}{*}{L01 cedolino singolo} & latenza (ms)      & 2.6 & 2.6 & 2.7 & 2.7 \\
 & throughput (op/s)  & 385 & 383 & 367 & 374 \\
 & CPU media/shard (\%) & 6 & 5 & 5 & 6 \\
 & disco letto (MB)   & 0.0 & 0.0 & 0.0 & 0.0 \\
 & disco scritto (MB) & 139.5 & 2.5 & 3.7 & 748.9 \\ \midrule
\multirow{5}{*}{L02 cedolini ditta/periodo} & latenza (ms)      & 3.0 & 3.1 & 3.2 & 3.1 \\
 & throughput (op/s)  & 329 & 325 & 316 & 321 \\
 & CPU media/shard (\%) & 6 & 6 & 6 & 6 \\
 & disco letto (MB)   & 0.0 & 0.0 & 0.0 & 0.0 \\
 & disco scritto (MB) & 5.6 & 153.9 & 112.1 & 4.1 \\ \midrule
\multirow{5}{*}{L03 costo per centro} & latenza (ms)      & 3.4 & 3.4 & 3.5 & 3.5 \\
 & throughput (op/s)  & 294 & 291 & 282 & 289 \\
 & CPU media/shard (\%) & 6 & 7 & 6 & 6 \\
 & disco letto (MB)   & 0.0 & 0.0 & 0.0 & 0.0 \\
 & disco scritto (MB) & 4.5 & 4.8 & 2.6 & 5.8 \\ \midrule
\multirow{5}{*}{L04 gender gap mediano} & latenza (ms)      & 26.0 & 61.6 & 122.2 & 185.2 \\
 & throughput (op/s)  & 38 & 16 & 8 & 5 \\
 & CPU media/shard (\%) & 14 & 15 & 18 & 22 \\
 & disco letto (MB)   & 0.0 & 0.0 & 0.0 & 0.0 \\
 & disco scritto (MB) & 5.4 & 6.0 & 5.2 & 2.9 \\ \midrule
\multirow{5}{*}{L05 top-N retribuzioni} & latenza (ms)      & 7.9 & 14.0 & 26.8 & 44.4 \\
 & throughput (op/s)  & 127 & 72 & 37 & 22 \\
 & CPU media/shard (\%) & 31 & 43 & 54 & 58 \\
 & disco letto (MB)   & 0.0 & 0.0 & 0.0 & 0.0 \\
 & disco scritto (MB) & 3.7 & 8.8 & 7.9 & 7.4 \\ \midrule
\multirow{5}{*}{L06 assenteismo per mese} & latenza (ms)      & 301.6 & 729.1 & 1684.9 & 2833.8 \\
 & throughput (op/s)  & 3 & 1 & 1 & 0 \\
 & CPU media/shard (\%) & 59 & 66 & 70 & 75 \\
 & disco letto (MB)   & 0.0 & 0.0 & 0.0 & 0.0 \\
 & disco scritto (MB) & 7.2 & 6.6 & 5.1 & 5.1 \\ \midrule
\multirow{5}{*}{L07 straordinari per centro} & latenza (ms)      & 3.6 & 3.7 & 3.8 & 3.8 \\
 & throughput (op/s)  & 274 & 270 & 261 & 265 \\
 & CPU media/shard (\%) & 8 & 8 & 8 & 8 \\
 & disco letto (MB)   & 0.0 & 0.0 & 0.0 & 0.0 \\
 & disco scritto (MB) & 3.3 & 3.9 & 4.4 & 2.6 \\ \midrule
\multirow{5}{*}{L08 trend costo lavoro} & latenza (ms)      & 29.0 & 68.8 & 149.4 & 255.5 \\
 & throughput (op/s)  & 34 & 14 & 7 & 4 \\
 & CPU media/shard (\%) & 49 & 58 & 64 & 69 \\
 & disco letto (MB)   & 0.0 & 0.0 & 0.0 & 0.0 \\
 & disco scritto (MB) & 9.9 & 4.6 & 5.4 & 10.8 \\ \midrule
\multirow{5}{*}{L09 organico per livello} & latenza (ms)      & 19.7 & 44.6 & 83.6 & 125.1 \\
 & throughput (op/s)  & 51 & 22 & 12 & 8 \\
 & CPU media/shard (\%) & 14 & 16 & 20 & 23 \\
 & disco letto (MB)   & 0.0 & 0.0 & 0.0 & 0.0 \\
 & disco scritto (MB) & 2.7 & 10.3 & 12.5 & 3.6 \\ \midrule
\multirow{5}{*}{L10 cedolino completo} & latenza (ms)      & 2.4 & 2.5 & 2.5 & 2.5 \\
 & throughput (op/s)  & 413 & 406 & 396 & 407 \\
 & CPU media/shard (\%) & 4 & 3 & 3 & 3 \\
 & disco letto (MB)   & 0.0 & 0.0 & 0.0 & 0.0 \\
 & disco scritto (MB) & 8.2 & 3.3 & 5.0 & 6.0 \\ \midrule
\multirow{5}{*}{L11 spesa per categoria} & latenza (ms)      & 34.2 & 80.0 & 178.0 & 308.0 \\
 & throughput (op/s)  & 29 & 12 & 6 & 3 \\
 & CPU media/shard (\%) & 48 & 56 & 61 & 65 \\
 & disco letto (MB)   & 0.0 & 0.0 & 0.0 & 0.0 \\
 & disco scritto (MB) & 3.1 & 4.8 & 4.3 & 3.5 \\ \midrule
\multirow{5}{*}{L12 ferie residue} & latenza (ms)      & 476.2 & 1174.1 & 2759.5 & 4786.3 \\
 & throughput (op/s)  & 2 & 1 & 0 & 0 \\
 & CPU media/shard (\%) & 57 & 62 & 68 & 89 \\
 & disco letto (MB)   & 0.0 & 0.0 & 0.0 & 0.0 \\
 & disco scritto (MB) & 12.5 & 5.7 & 6.1 & 12.3 \\ \midrule
\multirow{5}{*}{L13 riepilogo assenze} & latenza (ms)      & 33.8 & 34.4 & 42.8 & 38.6 \\
 & throughput (op/s)  & 30 & 29 & 23 & 26 \\
 & CPU media/shard (\%) & 17 & 17 & 17 & 21 \\
 & disco letto (MB)   & 0.0 & 0.0 & 0.0 & 0.0 \\
 & disco scritto (MB) & 3.1 & 10.4 & 9.2 & 2.9 \\ \midrule
\bottomrule
\end{longtable}
\endgroup


\begin{figure}[htbp]
  \centering
  \includegraphics[width=0.82\textwidth]{fig-scaling-dati-mongo}
  \caption{MongoDB, scaling per dati (tutte le 13 letture, 4 shard fissi, scala log): L06 e L12 crescono a
    dismisura col volume, le single-shard restano piatte.}
  \label{fig:mongo-dati}
\end{figure}

\begin{lstlisting}[style=sobrio, language=bash, caption={Esecuzione dello scaling per dati (MongoDB): 4 shard fissi, lotti crescenti di tenant, misura a caldo.}]
export MONGO_PASSWORD=...
bash scala-dati-mongo.sh 64 200 400 800     # +200/+400/+800 tenant -> 230, 630, 1430 ditte
\end{lstlisting}

\subsection{Scritture (MongoDB)}
\label{sec:mongo-scritture}
Le scritture (4 shard, N${=}300$, tenant fisso) confermano la co-locazione e, su U03, producono il contrasto
pro-relazionale più marcato dell'intero lavoro (\Cref{tab:mongo-scritture}). Le scritture \emph{intra-tenant}
(I01, I03, U01, U02, U04) colpiscono \textbf{un solo shard} (quello del tenant): il carico è tutto su
\texttt{10.0.1.4} (CPU ${\sim}5\,\%$), gli altri tre inattivi, esattamente come in Citus. \textbf{U03}
invece, che aggiorna una categoria \emph{denormalizzata} in tutti i cedolini, è un \texttt{updateMany} che
tocca \textbf{tutti e 4 gli shard al 91--99\,\% di CPU}: 356\,ms per operazione e appena 2.8\,op/s. In Citus
la stessa modifica è \textbf{una riga} su una reference table (21\,ms, 47\,op/s): un fattore ${\sim}17\times$
a favore del relazionale.

\begin{table}[htbp]
  \centering
  \caption{MongoDB, scritture a 4 shard (N${=}300$): la co-locazione (una scrittura, uno shard) e il mass-update U03.}
  \label{tab:mongo-scritture}
  \begin{tabular}{llrrl}
    \toprule
    \textbf{Op.} & \textbf{Descrizione} & \textbf{lat (ms/op)} & \textbf{tps} & \textbf{dove scrive} \\
    \midrule
    U02 & licenziamento                & 3.6   & 279 & solo 1 shard \\
    U01 & coordinate bancarie          & 3.6   & 279 & solo 1 shard \\
    U04 & update concorrente           & 3.6   & 277 & solo 1 shard \\
    I01 & inserimento timbratura       & 8.6   & 116 & solo 1 shard \\
    I03 & elaborazione cedolino        & 11.1  & 90  & solo 1 shard \\
    U03 & modifica dato di riferimento & 355.9 & 2.8 & \textbf{tutti i 4 shard (91--99\,\% CPU)} \\
    \bottomrule
  \end{tabular}
\end{table}

\begin{figure}[htbp]
  \centering
  \includegraphics[width=0.85\textwidth]{fig-scritture}
  \caption{Scritture a confronto (4 nodi/shard, N${=}300$): latenza per operazione, Citus vs MongoDB.}
  \label{fig:scritture}
\end{figure}

\begin{lstlisting}[style=sobrio, language=bash, caption={Esecuzione delle scritture (MongoDB): N${=}300$ ripetizioni per query, tenant fisso.}]
bash esegui-scritture-mongo.sh 300 queries/document/scritture/*.js
\end{lstlisting}

\subsection{Comportamento in presenza di guasti (MongoDB)}
\label{sec:mongo-guasti}
Con lo stesso protocollo del test di guasto Citus, si esegue L04 in loop mentre lo shard che ospita i dati
del tenant (\texttt{10.0.1.4}) viene spento e riacceso (10 run). Il comportamento è \textbf{CP}: mentre lo
shard è giù la query fallisce e le righe crollano da 59\,299 a \textbf{0} (nessun risultato parziale), per
tornare al valore pieno al rientro; il recupero è automatico (${\sim}0.25$\,s dopo il riavvio) e
l'indisponibilità dura esattamente il tempo di down (${\sim}15$\,s) (\Cref{tab:mongo-guasti}).

\begin{table}[htbp]
  \centering
  \caption{MongoDB, guasto di uno shard sotto carico di lettura (10 run).}
  \label{tab:mongo-guasti}
  \begin{tabular}{lrrll}
    \toprule
    \textbf{Shard spento} & \textbf{run} & \textbf{error-rate} & \textbf{finestra / recupero} & \textbf{righe (su$\to$giù$\to$su)} \\
    \midrule
    10.0.1.4 (con i dati) & 10 & 7.0\,\% & 10.1$\to$25.2\,s / ${\sim}0.25$\,s & 59\,299 $\to$ \textbf{0} $\to$ 59\,299 \\
    \bottomrule
  \end{tabular}
\end{table}

\begin{figure}[htbp]
  \centering
  \includegraphics[width=0.8\textwidth]{fig-guasti}
  \caption{Guasto di un nodo sotto carico L04, Citus e MongoDB: righe restituite (normalizzate) nel tempo. In
    entrambi crollano a 0 durante il down (banda) e recuperano al rientro (comportamento CP).}
  \label{fig:guasti}
\end{figure}

Un dettaglio metodologicamente importante: l'error-rate è solo ${\sim}7\,\%$, contro il ${\sim}51\,\%$ di
Citus, \emph{pur avendo la stessa indisponibilità reale} (15\,s). La ragione è il \textbf{diverso
comportamento del client}: il \texttt{mongos} \textbf{si blocca in attesa} dello shard irraggiungibile fino
al timeout (5\,s per operazione), quindi nella finestra di down passano poche operazioni-errore ``lunghe'';
Citus invece \textbf{rifiuta la connessione all'istante} (fail-fast) e ne accumula molte di più. L'error-rate
grezzo non è dunque direttamente comparabile: va letto insieme al \emph{modo} di fallire, cioè \textbf{timeout
bloccanti} in Mongo contro \textbf{errori rapidi} in Citus.

\begin{lstlisting}[style=sobrio, language=bash, caption={Esecuzione del test di guasto (MongoDB): L04 in loop, arresto/riavvio dello shard con i dati.}]
bash test-guasto-mongo.sh queries/document/letture/L04_gender_gap_mediano.js 10.0.1.4 10 15 45
\end{lstlisting}

\section{Confronto tra i due sistemi}
\label{sec:confronto}

Il confronto si articola su più assi (\Cref{tab:confronto,tab:confronto-completo}):
\begin{itemize}
  \item \textbf{Analitica}: dipende dal tipo di query. Dove l'aggregato è servito dall'embedding (percentili
        L04, spesa per categoria L11) MongoDB è più veloce, perché evita le join (a 1 nodo/shard L04 7.2 vs
        71.9\,ms, L11 12.4 vs 75.6\,ms). Ma sulle query che \textbf{attraversano più entità} (assenteismo L06,
        ferie residue L12, riepilogo assenze L13, che toccano assenze e ratei) vince Citus, e di molto: a pari
        scala (${\sim}230$--$240$ ditte, 4 nodi) L06 51 vs 302\,ms, L12 28 vs 476\,ms, L13 3.6 vs 34\,ms.
        MongoDB paga lo scan degli array annidati sull'intera collezione, Citus sfrutta i join indicizzati
        sulle tabelle normalizzate.
  \item \textbf{Scaling}: Citus distribuisce e scala \emph{sempre} (rebalance esplicito, ${\sim}4\times$ con
        4 nodi); MongoDB invece dipende dalla scala: a dati piccoli il balancer non distribuisce (nessuno
        scaling), a volume grande sì (tutti gli shard attivi, cfr.\ \Cref{sec:mongo-dati}).
  \item \textbf{Scritture}: entrambi co-locano la scrittura di un tenant su un solo nodo/shard, ma sul dato
        di riferimento (U03) il relazionale è nettamente più efficiente (1 riga, 21\,ms vs mass-update,
        356\,ms).
  \item \textbf{Guasti}: entrambi CP con la stessa indisponibilità reale, ma \emph{modo} di fallire opposto
        (fail-fast vs timeout bloccante).
  \item \textbf{Query mirate} (cedolino singolo): comparabili (${\sim}$1--2\,ms) su entrambi.
\end{itemize}

\begin{table}[htbp]
  \centering
  \caption{Confronto in latenza (ms) a pari numero di nodi (1) e capacità di scalare.}
  \label{tab:confronto}
  \begin{tabular}{lrrl}
    \toprule
    \textbf{Query} & \textbf{Citus} & \textbf{MongoDB} & \textbf{Scaling coi nodi} \\
    \midrule
    L04 gender gap    & 71.9 & \textbf{7.2}  & Citus ${\sim}4\times$; Mongo no (scala piccola) \\
    L11 spesa/categoria & 75.6 & \textbf{12.4} & Citus ${\sim}4\times$; Mongo no (scala piccola) \\
    L01 cedolino singolo & \textbf{1.5} & 2.3 & entrambi piatti (mirata) \\
    \bottomrule
  \end{tabular}
\end{table}

\begin{figure}[htbp]
  \centering
  \begin{minipage}[t]{0.49\textwidth}
    \centering
    \includegraphics[width=\linewidth]{fig-scaling-nodi-confronto}
    \subcaption{Scaling per nodi su L04: Citus scala, Mongo piatto (scala piccola).}
    \label{fig:scaling-confronto}
  \end{minipage}\hfill
  \begin{minipage}[t]{0.49\textwidth}
    \centering
    \includegraphics[width=\linewidth]{fig-letture-assolute}
    \subcaption{Latenza assoluta delle 13 letture a scala base (log).}
    \label{fig:letture-assolute}
  \end{minipage}
  \caption{Confronto Citus vs MongoDB: capacità di scalare coi nodi (a) e velocità assoluta in lettura (b).}
  \label{fig:confronto}
\end{figure}

\begin{table}[htbp]
  \centering
  \caption{Confronto complessivo Citus vs MongoDB sui diversi assi misurati.}
  \label{tab:confronto-completo}
  \begin{tabularx}{\textwidth}{@{}lXX@{}}
    \toprule
    \textbf{Aspetto} & \textbf{PostgreSQL + Citus} & \textbf{MongoDB} \\
    \midrule
    Distribuzione dei dati            & esplicita (\texttt{rebalance}, deterministica) & automatica (balancer, soggetta a soglia) \\
    Scaling per nodi (dati piccoli)   & ${\sim}4\times$ con 4 nodi (cross-shard)       & nessuno (balancer sotto soglia) \\
    Scaling per dati (volume grande)  & cross-shard super-lineari                     & idem, e il balancer distribuisce \\
    Analitica sull'aggregato (L04, L11) & più lenta (join)                           & più veloce (embedding) \\
    Query multi-entità (L06, L12, L13)  & \textbf{più veloce} (join indicizzati)     & molto più lenta (array annidati) \\
    Scrittura intra-tenant            & 1 nodo (co-locazione)                         & 1 shard (co-locazione) \\
    Dato di riferimento (U03)         & \textbf{1 riga: 21\,ms, 47\,op/s}             & mass-update: 356\,ms, 2.8\,op/s \\
    Guasto (modo di fallire)          & fail-fast (err ${\sim}51\,\%$)                & timeout bloccante (err ${\sim}7\,\%$) \\
    Guasto (indisponibilità reale)    & ${\sim}15$\,s, recupero ${\sim}0.5$\,s        & ${\sim}15$\,s, recupero ${\sim}0.25$\,s \\
    \bottomrule
  \end{tabularx}
\end{table}

Il confronto diretto sulle \textbf{scritture} (\Cref{tab:confronto-scritture}) dà esiti opposti a seconda
dell'operazione: sull'inserimento del cedolino (I03) è molto più rapido MongoDB (11 vs 57\,ms), perché scrive
un solo documento invece di una transazione su sette tabelle; sull'aggiornamento del dato di riferimento (U03)
è molto più rapido Citus (21 vs 356\,ms), perché tocca una riga su una reference table invece di aggiornare in massa tutti gli
shard. Le scritture semplici (U01/U02/U04) sono equivalenti. Sull'\textbf{analitica}
(\Cref{tab:confronto}) MongoDB è più veloce in assoluto grazie all'embedding.

\begin{table}[htbp]
  \centering
  \caption{Scritture a confronto (4 nodi/shard, N${=}300$): latenza e throughput per operazione.}
  \label{tab:confronto-scritture}
  \begin{tabular}{llrrrr}
    \toprule
    & & \multicolumn{2}{c}{\textbf{Citus}} & \multicolumn{2}{c}{\textbf{MongoDB}} \\
    \cmidrule(lr){3-4}\cmidrule(lr){5-6}
    \textbf{Op.} & \textbf{Descrizione} & \textbf{ms/op} & \textbf{op/s} & \textbf{ms/op} & \textbf{op/s} \\
    \midrule
    U02 & licenziamento          & 1.7  & 605 & 3.6   & 279 \\
    U01 & coordinate bancarie    & 3.7  & 273 & 3.6   & 279 \\
    U04 & update concorrente     & 3.8  & 266 & 3.6   & 277 \\
    I01 & inserimento timbratura & 10.6 & 94  & 8.6   & 116 \\
    I03 & elaborazione cedolino  & 57.4 & 17  & \textbf{11.1}  & \textbf{90} \\
    U03 & dato di riferimento    & \textbf{21.1} & \textbf{47} & 355.9 & 2.8 \\
    \bottomrule
  \end{tabular}
\end{table}

Il contrasto tra i modelli si vede meglio sulle coppie deliberate. Il \textbf{cedolino completo} è
pro-documentale: una \texttt{findOne} su un documento (\Cref{lst:l10-mongo}) contro l'assemblaggio da sette
tabelle in Citus. Il \textbf{dato di riferimento} è pro-relazionale: una riga su una reference table
(\Cref{lst:u03-citus}) contro un \texttt{updateMany} di centinaia di documenti denormalizzati in Mongo.

\begin{lstlisting}[style=sobrio, language=Java, caption={Cedolino completo in MongoDB: un solo documento (pro-documentale).}, label={lst:l10-mongo}]
db.cedolini.findOne({ tenant_id: 1, dipendente_id: 7, anno: 2026, mese: 6 });
\end{lstlisting}

\begin{lstlisting}[style=sobrio, language=SQL, caption={Modifica di un dato di riferimento in Citus: una riga (pro-relazionale).}, label={lst:u03-citus}]
UPDATE tipo_voce SET categoria = 'WELFARE' WHERE codice = '0401';   -- 1 riga (in Mongo: updateMany su N documenti)
\end{lstlisting}
